\documentclass[a4paper,12pt]{book}

\usepackage[unicode,colorlinks=true,linkcolor=blue]{hyperref}
% здесь подключении шрифтов в русскими буквами
\usepackage[T2A]{fontenc}
\usepackage[utf8]{inputenc}
\usepackage[english,russian]{babel}
\usepackage{amsmath,amsthm,amssymb,amsfonts,mathtext,cite,enumerate,float}
%\usepackage[dvips]{graphicx}
\usepackage[pdftex]{graphicx}
\graphicspath{{images/}}
\usepackage{fix-cm}

\makeatletter
\renewcommand{\@biblabel}[1]{#1.}
\makeatother

\usepackage{geometry}  % Меняем поля страницы
\geometry{left=2cm}    % левое поле
\geometry{right=1.5cm} % правое поле
\geometry{top=1cm}     % верхнее поле
\geometry{bottom=2cm}  % нижнее поле

\newcommand{\HRule}{\rule{\linewidth}{0.5mm}}

\renewcommand{\theenumi}{\arabic{enumi}}                                       % Меняем везде перечисления на цифра.цифра
\renewcommand{\labelenumi}{\arabic{enumi}.}                                     % Меняем везде перечисления на цифра.цифра
\renewcommand{\theenumii}{.\arabic{enumii}}                                    % Меняем везде перечисления на цифра.цифра
\renewcommand{\labelenumii}{\arabic{enumi}.\arabic{enumii}.}                   % Меняем везде перечисления на цифра.цифра
\renewcommand{\theenumiii}{.\arabic{enumiii}}                                  % Меняем везде перечисления на цифра.цифра
\renewcommand{\labelenumiii}{\arabic{enumi}.\arabic{enumii}.\arabic{enumiii}.} % Меняем везде перечисления на цифра.цифра

\theoremstyle{plain}
\newtheorem{theorem}{Теорема}[section]
\newtheorem*{consequence}{Следствие}
\newtheorem{lemma}[theorem]{Лемма}
\theoremstyle{definition}
\newtheorem{definition}{Определение}[section]
\theoremstyle{remark}
\newtheorem*{comment}{Комментарий}
\newtheorem{example}{Пример}[section]
\newtheorem{remark}{Замечание}[section]


\begin{document}
  \begin{titlepage}
  \begin{center}

    % Upper part of the page. The '~' is needed because \\
    % only works if a paragraph has started.
    %\includegraphics[width=0.15\textwidth]{./logo}~\\[1cm]

    \textsc{\LARGE НИЯУ МИФИ}\\[1.5cm]

    \textsc{\Large 3 семестр факультет КиБ}\\[0.5cm]

    % Title
    \HRule \\[0.4cm]
    {\huge \bfseries Математический анализ\\[0.4cm]}

    \HRule \\[1.5cm]

    \begin{minipage}{0.4\textwidth}
      \begin{flushleft} \large
        \emph{Автор:}\\
        Тропин \textsc{А.Г.}
      \end{flushleft}
    \end{minipage}
    \begin{minipage}{0.4\textwidth}
      \begin{flushright} \large
        \emph{Лектор:} \\
        Теляковский \textsc{Д.С.}
      \end{flushright}
    \end{minipage}

    \vfill
    \begin{flushleft}
      e-mail: \href{mailto:andrewtropin@gmail.com}{andrewtropin@gmail.com} \\
      github: \href{http://github.com/abcdw/mephi}{abcdw/mephi}
    \end{flushleft}
    {\large \today}
  \end{center}
\end{titlepage}

  \tableofcontents

  \part{Функциональный последовательности и ряды}
  \chapter{Числовые ряды}
  \section{Определение}

\begin{definition}
  $U_1 + U_2 + U_3 + \dots = \sum\limits_{k = 0}^{\infty}{U_k}$
\end{definition}

\begin{definition}[Частичная сумма]
  $S_n = \sum\limits_{k = 0}^n U_k$
\end{definition}

\begin{definition}
  Ряд сходится, если $\exists \lim\limits_{n \to \infty}
                      \sum\limits_0^\infty U_k = S$
\end{definition}

\begin{definition}
  $\{a_n\} = a_0 + \sum\limits_1^n (a_k - a_{k-1})$, где $(a_k - a_{k - 1}) = U_k$
\end{definition}

\begin{theorem}[Критерий Коши]
  \label{th111}
  Ряд сходится, тогда и только тогда, когда он удовлетворяет условию Коши: \\
  $$\forall \varepsilon > 0, \ \exists N~=~N(\varepsilon), \
  \forall n \geq N, \forall p: \ |\sum\limits_{k=n+1}^{n+p} U_k|~=~|U_{n+1}
  + \dots + U_{n+p}|~=~|S_{n+p} - S_{n}| < \varepsilon$$
\end{theorem}

\begin{proof}
  $\sum U_k$ - сходится $\Leftrightarrow \{S_n\}$
  $$\forall \varepsilon > 0, \ \exists N: \forall n \geq N, \forall p: \
  |S_{n+p} - S_{n}| < \varepsilon$$
\end{proof}
\begin{consequence}[Необходимое условие сходимости]
  Если $\sum U_k$ сходится, то $U_k \to 0$, при $k~\to~\infty$
\end{consequence}
\begin{proof}
  Если $\sum U_k$ сходится, то выполняется Критерий Коши. При $p = 1$
  $$\forall \varepsilon > 0, \exists N, \forall n \geq N:
  |S_{n+1} - S_n|~=~|U_{n+1}| < \varepsilon$$
\end{proof}
\begin{consequence}
  Отбрасывание или добавление любого конечного числа членов ряда на его
  сходимость не влияет.
\end{consequence}

\begin{example}
  $\sum\limits_0^\infty z^n,$
  $S_n(z) = \sum\limits_0^n z_n = \frac{1-z^{n+1}}{1-z}$.
  При $n \to \infty, \ S_n(z) = \frac{1}{1 - z}, \ |z| < 1$. \\
  $S_n(z)$ не имеет придела при $|z| \geq 1$.
\end{example}

\section{Действия с рядами}

\begin{theorem}
  Ряды $\sum U_k$ и $\sum V_k$ сходятся, $\alpha$ --- комплексное число, тогда
  \begin{align}
    \label{series:property1}
    \sum\alpha U_k &= \alpha\sum U_k \\
    \label{series:property2}
    \sum (U_k \pm V_k) &= \sum U_k \pm \sum V_k
  \end{align}
\end{theorem}
\begin{proof}
  Доказательство свойства (\eqref{series:property1}):
  $$\sum\limits_{k=0}^\infty \alpha U_k = \lim\limits_{n \to \infty}
  \sum\limits_0^n \alpha U_k = \alpha \lim\limits_{n \to \infty} \sum\limits_0^n U_k
  = \alpha \sum\limits_{k = 0}^\infty U_k$$
\end{proof}
\begin{proof}
  Доказательство свойства (\eqref{series:property2}):
  \begin{gather*}
      \sum\limits_0^\infty U_k \pm \sum\limits_{0}^{\infty} V_k
      = \lim\limits_{n \to \infty} \sum\limits_{0}^{n} U_k
      \pm \lim\limits_{n \to \infty} \sum\limits_{0}^{n} V_k = \\
      = \lim\limits_{n \to \infty} \sum\limits_{0}^{n} U_k
      \pm \sum\limits_{0}^{n} V_k
      = \lim\limits_{n \to \infty} \sum\limits_{0}^{n} (U_k \pm V_k)
      = \sum\limits_{0}^{n} (U_k \pm V_k)
  \end{gather*}
\end{proof}

\begin{remark}
  Из сходимости $\sum (U_k \pm V_k) \not \Rightarrow$ сходимость
  $\sum U_k$ и $\sum V_k$
\end{remark}
\begin{remark}
  Если $\sum U_k$ сходится, то можно группировать, не меняя порядка.
\end{remark}
\begin{example}
  \begin{gather*}
    \sum (1 - 1) \\
    (1 - 1) + (1 - 1) + \dots \\
    1 - (1 - 1) - (1 - 1) \dots
  \end{gather*}
\end{example}
\begin{comment}
  Нельзя раскрывать скобки и переставлять члены.
\end{comment}
\subsection{Ряды с неотрицательными членами}
$U_k \geq 0$, $S_n = \sum\limits_{0}^{n} U_k$ - не убывающая последовательность. \\
$\sum\limits_{0}^{n} U_k$ - сходится $\Leftrightarrow$ $\{S_n\} $ --- ограничена
\begin{comment}
  Сходимость ряда эквивалентна ограниченности $S_n$
\end{comment}

\begin{theorem}
    $$U_k \geq 0, V_k \geq 0, \ \forall k:$$
    \begin{enumerate}
      \item
        Если $0 \leq U_k \leq V_k$, то если $\sum V_k$ сходится $\Rightarrow$
        $\sum U_k$ сходится и если
        $\sum U_k$ расходится $\Rightarrow$ $\sum V_k$ расходится.
      \item
        Если $\lim\limits_{n\to \infty} \frac{U_k}{V_k} = A > 0$, то ряды
        сходятся или расходятся одновременно.
    \end{enumerate}
\end{theorem}

\pagebreak
\begin{proof}
  \hfill
  \begin{enumerate}
    \item
      $\forall n$ верно неравенство $0 \leq \sum\limits_{0}^{n} U_k \leq
      \sum\limits_{0}^{n} V_k$
    \item
      $\forall \varepsilon > 0 \ \varepsilon < A \ \exists N : \forall n \geq N
      \Rightarrow 0 < A - \varepsilon < \frac{U_k}{V_k} < A + \varepsilon$ \\
      $0 < (A - \varepsilon) \cdot V_k < U_k < (A + \varepsilon) \cdot V_k$ \\
      Пусть $U_k$ --- сходится, тогда из доказанного выше 1ого пункта следует
      $(A - \varepsilon) \cdot V_k$ --- сходится
      $\Rightarrow \sum V_k$ сходится
      $\Rightarrow \sum (A + \varepsilon) \cdot V_k$ сходится
      $\Rightarrow \sum U_k$ сходится.
  \end{enumerate}
\end{proof}

\begin{remark}
  Вместо существования предела $\lim\limits_{n \to \infty} \frac{U_k}{V_k}$
  достаточно предположить, что существуют такие числа
  p и q > 0, такие что $0 < q < \frac{U_k}{V_k} < p, \ \forall k$
\end{remark}

\begin{theorem}[Признак Даламбера]
  $$\sum U_k, \ U_k > 0$$
  \begin{enumerate}
    \item Если $\exists q$ такое что:
      $\ \forall k \ \frac{U_{k+1}}{U_k} \leq q < 1$ сходится
    \item Если $\exists \lim\limits_{k \to \infty} \frac{U_{k+1}}{U_k} = q$, то:
    \begin{itemize}
      \item при $q < 1$ сходимость
      \item при $q > 1$ расходимость
      \item при $q = 1$ неизвестно (нужно провести дополнительные исследования)
    \end{itemize}
  \end{enumerate}
\end{theorem}

\begin{proof}
  Идея докозательства - сравнение с геометрической прогрессией.
  \begin{enumerate}
    \item $k = 0, 1, \dots, n$;
      $U_k = U_0 \cdot \frac{U_1}{U_0} \frac{U_2}{U_1} \cdots \frac{U_k}{U_{k-1}}
      < U_0 \cdot q^k$
      \begin{comment}
        $\frac{U_k}{U_{k-1}} < q, \ \forall k$
      \end{comment}
      $q < 1$, тогда $\sum U_0 \cdot q^k$ --- сходящаяся геометрическая прогрессия. \\
      $U_k = U_0 \cdot \frac{U_1}{U_0}\frac{U_2}{U_1}\cdots\frac{U_k}{U_{k-1}}
      \geq U_0 > 0$
      \begin{comment}
        $\frac{U_k}{U_{k-1}} \geq 1, \ \forall k$
      \end{comment}
      $U_k \not \to 0 \Rightarrow$ не выполняется необходимое условие сходимости.
    \item Пусть $\lim\limits_{k \to \infty} \frac{U_{k+1}}{U_k} = q$ \\
      $\forall \varepsilon > 0, \ \exists K : \ \forall k \geq K$
      выполняется неравенство
      $q - \varepsilon < \frac{U_{k+1}}{U_k} < q + \varepsilon$
      \begin{itemize}
        \item Если $q < 1$,выберем такое $\varepsilon$, что 
          $q + \varepsilon < 1$, для
          $\forall k \geq K(\varepsilon)$. \\
          $\frac{U_{k+1}}{U_k} < q + \varepsilon < 1 \Rightarrow$
          сходится по первой части.
        \item Если $q > 1$, то выберем $\varepsilon$ так, чтобы
          $q - \varepsilon > 1$, для
          $\forall k \geq K(\varepsilon)$. \\
          $\frac{U_{k+1}}{U_k} > q - \varepsilon > 1, \ \Rightarrow$
          расходится по первой части.
      \end{itemize}
    \end{enumerate}
\end{proof}

\pagebreak
\begin{theorem}[Признак Коши]
  $$\sum U_k, U_k \geq 0$$
  \begin{enumerate}
    \item Если $\exists q < 1$ и $\forall k > K : $ выполняется
      $\sqrt[k]U_k \leq q < 1$, то ряд сходится, а если
      $\forall k \ \sqrt[k]U_k~\geq~1$, то расходится.
    \item Если $\exists \lim\limits_{k \to \infty} \sqrt[k]U_k = q, (q \geq 0)$,
      то
      \begin{itemize}
        \item $q < 1$ - сходится
        \item $q > 1$ - расходится
        \item $q = 1$ - нужны дополнительные исследования
      \end{itemize}
  \end{enumerate}
\end{theorem}
\begin{remark}
  $\lim\limits_{n \to \infty} \sqrt[n]U_n$ можно рассматривать вместо
  $\overline {\lim\limits_{k \to \infty}} \sqrt[k]U_k$
\end{remark}

\begin{proof}
  Сравнение с геометрической прогрессией
  \begin{enumerate}
    \item Если $\forall k\ \sqrt[k]U_k \leq q < 1 \Rightarrow U_k \leq q^k$
      --- сходящаяся геометрическая прогрессия. \\
      Если $\forall k \sqrt[k]U_k \geq 1 \Rightarrow U_k \geq 1$
      --- не выполняется необходимое условие сходимости.
    \item Если $\lim\limits_{k \to \infty} \sqrt[k]U_k = q$, то
      $\forall \varepsilon > 0 \ \exists K = K(\varepsilon) :
      \forall k \geq K, (q - \varepsilon) < \sqrt[k]U_k < (q + \varepsilon)$\\
      $(q - \varepsilon)^k < U_k < (q + \varepsilon)^k$
      \begin{itemize}
        \item При $q < 1$ выберем $\varepsilon$ так, чтобы $q + \varepsilon < 1$,
          тогда $U_k < (q + \varepsilon)^k < 1$
          --- сходящаяся геометрическая прогрессия.
        \item При $q > 1$ выберем $\varepsilon$ так, чтобы $q - \varepsilon > 1$,
          тогда $U_k > (q - \varepsilon)^k > 1$
          --- не выполняется необходимое условие сходимости.
      \end{itemize}
  \end{enumerate}
\end{proof}

\begin{definition}
  Дана $\{a_n\}$ и пусть $\overline {\lim\limits_{n\to \infty}} a_n$
  --- наибольший из частичных пределов, тогда:
  $$\forall \{a\} \ \exists \overline{\lim\limits_{n \to \infty}} a_n = A \ or \ \infty $$
  \begin{comment}
    $A$ --- число.
  \end{comment}
  \begin{itemize}
    \item Если $\overline{\lim\limits_{n \to \infty}} a_n = +\infty \Rightarrow
      \{a_n\}$ --- неограничена сверху $\Rightarrow
      \overline{\lim\limits_{k \to \infty}} \sqrt[k]U_k = +\infty$
      --- неограничена сверху. $U_k$ неограничена сверху и не выполняется
      необходимое условие.
    \item Если $\overline{\lim\limits_{n \to \infty}} a_n = A$, тогда
      $\forall \varepsilon \in (A - \varepsilon, A + \varepsilon)$ бесконечно
      много членов $\{a_n\}$: \\
      \begin{itemize}
        \item $\overline{\lim\limits_{k \to \infty}} \sqrt[k]U_k = q < 1$.
          Выберем $\varepsilon$ так, чтобы $q + \varepsilon < 1 \Rightarrow$
          $\exists K: \forall k \geq K, \ \sqrt[k]U_k < q + \varepsilon < 1$ по признаку
          Коши. \\
        \item $\overline{\lim\limits_{k \to \infty}} \sqrt[k]U_k = q > 1$.
          Выберем $\varepsilon$ так, чтобы $q - \varepsilon > 1 \Rightarrow$
          $\forall K \ \exists k \geq K: \sqrt[k]U_k~>~q~-~\varepsilon~>~1 \\
          \Rightarrow~U_k~>~1$
      \end{itemize}
  \end{itemize}
\end{definition}

\section{Интегральный признак сходимости рядов с неотрицательными членами}
\begin{theorem}
  Если $f(x)$ не отрицательна и убывает на $x \geq 1$, то ряд
  \begin{gather}
    \label{th131:series1}
    \sum\limits_{n = 1}^{\infty} f(n) < \infty
  \end{gather}
  сходится тогда и только тогда, когда сходится интеграл:
  \begin{gather}
    \label{th131:integral1}
    \int\limits_1^{+\infty} f(x) dx
  \end{gather}
  то есть $ \int\limits_1^{+\infty} f(x) dx < \infty$. \\
  \begin{itemize}
    \item $\sum a_n < \infty$ --- сходится
    \item $\sum a_n = \infty$ --- расходится
  \end{itemize}
\end{theorem}

\begin{proof}
  Если $k \leq x \leq k + 1, \ k = 1, 2, \dots,$ то, в силу убывания функции
  получаем неравенство:
  $$f(k) \geq f(x) \geq f(k + 1)$$
  Интегрируя по отрезку $[k, k + 1]$ получим:
  $$f(k) \geq \int\limits_k^{k+1} f(x) dx \geq f(k+1), \ k = 1, 2, \dots \ .$$

  \begin{gather}
    \label{th131:series2}
    \sum\limits_{k = 1}^{n} f(k+1) \leq \int\limits_1^{n+1} f(x) dx \leq
    \sum\limits_{k = 1}^{n} f(k)
  \end{gather}
  Пусть $S_n = \sum\limits_{k = 1}^{n} f(k),$ тогда \eqref{th131:series2} примет вид:
  \begin{gather}
    \label{th131:series3}
    S_n + 1 - f(1) \leq \int\limits_1^{n+1} f(x) dx \leq S_n
  \end{gather}
  Если ряд \eqref{th131:series1} сходится и его сумма равна $S$, то $S_n \leq S$,
  и $\int\limits_1^{n+1} f(x) dx \leq S, \ \forall n \in \mathbb{N}$. \\
  $\forall b > 1, n + 1 > b$ имеем:
  $$\int\limits_1^b f(x) dx \leq \int\limits_1^{n+1} f(x) dx \leq S$$
  В силу неотрицательности функции $f(x)$ интеграл сходится. \\
  Пусть наоборот, интеграл \eqref{th131:integral1} сходится, тогда из
  \eqref{th131:series3} следует:
  $$S_{n+1} \leq f(1) + \int\limits_1^{n+1} f(x) dx
  \leq f(1) + \int\limits_1^\infty f(x) dx$$
  Тем самым, последовательность сумм $\{S_n\}$ ряда \eqref{th131:series1}
  ограничена сверху, и поэтому этот ряд сходится.
\end{proof}

\begin{example}
  \label{th131:example131}
  \begin{gather}
    \label{ex131:series1}
    \sum\limits_{n = 1}^{\infty} \frac{1}{n^\alpha}, \alpha \in \mathbb{R}
  \end{gather}
  Положим $f(x) = \frac{1}{x^\alpha}$, тогда $f(n) = \frac{1}{n^\alpha}$ \\
  Поскольку $\int\limits_1^{+\infty} \frac{dx}{x^\alpha}$:
  \begin{itemize}
    \item При $\alpha > 1$ сходится
    \item При $\alpha \leq 1$ расходится
  \end{itemize}
  Тогда ряд \eqref{ex131:series1} сходится тогда и только тогда, когда $\alpha > 1$.
  При $\alpha < 0 \ $ дробь $ \ \frac{1}{n^\alpha} \geq 1$.
\end{example}

\section{Признак сходимости для чередующихся рядов}
Рассмотрим ряды с действительными числами, которые то положительные, то
отрицательные.

\begin{theorem}[Лейбница]
  Если
  \begin{gather}
    \lim\limits_{n \to \infty} U_n = 0 \label{th141:lim1} \\
    U_n \geq U_{n+1} > 0, \ n = 1, 2, \dots \label{th141:ineq1}
  \end{gather}
  то знакочередеющийся ряд
  \begin{gather}
    \sum\limits_{n = 1}^{+\infty} {(-1)}^{n+1} U_n \label{th141:sum1}
  \end{gather}
  сходится, при этом если $S$ --- сумма ряда, а $S_n$ --- его $n$-ая частичная
  сумма,\\ то $\forall n: n = 1, 2, \dots$
  \begin{gather}
    |S - S_n| \leq U_{n+1} \label{th141:ineq2}
  \end{gather}
\end{theorem}

\begin{proof}
  Заметим, что частичная суммы $S_n$ с четными номерами возрастают:
  $$S_{2k} = (U_1 - U_2) + (U_3 - U_4) + \dots + (U_{2k - 1} - U_{2k}), \
  k = 1, 2, \dots $$
  Так что выполняется неравенство $S_{2k + 2} \geq S_k$. Кроме того, они
  ограничены сверху:
  $$S_{2k} = U_1 - (U_2 - U_3) - \dots - (U_{2k - 2} - U_{2k - 1}) - U_{2k}, \
  S_{2k} < U_1$$
  Поэтому последовательность $\{S_{2k}\}$ сходится
  \begin{gather}
    \lim\limits_{k \to \infty} S_{2k} = S \label{th141:lim2}
  \end{gather}
  Поскольку $S_{2k + 1} = S_{2k} + U_{2k + 1}$ и $U_{2k + 1} \to 0$ при
  $k \to \infty$, то
  \begin{gather}
    \lim\limits_{k \to \infty} S_{2k + 1} = S \label{th141:lim3}
  \end{gather}
  Из \eqref{th141:lim2} и \eqref{th141:lim3} следует, что
  $\lim\limits_{n \to \infty} S_n = S$ \\
  При этом, нетрудно увидеть, что
  \begin{gather}
    S_{2k} \leq S \leq S_{2k + 1} \leq S_{2k - 1}, \ \forall k \label{th141:ineq3}
  \end{gather}
  Из неравенства \eqref{th141:ineq3} следует, что
  \begin{align*}
    S - S_{2k} \leq S_{2k + 1} - S_{2k} &= U_{2k + 1} \\
    S_{2k - 1} - S \leq S_{2k - 1} - S_{2k} &= U_{2k}, \ k = 1, 2, \dots
  \end{align*}
  Это и означает, что $\forall n \in \mathbb{N}$ выполняется
  неравенство \eqref{th141:ineq2}.
\end{proof}

\section{Преобразование Абеля}
\begin{theorem}
  Пусть $a_k \in \mathbb{C}$, $b_k \in \mathbb{C}$, $k = 1, \dots, n$;
  $B_k = b_1 + \dots + b_k$, тогда
  \begin{gather}
    \sum\limits_{k = 1}^{n} a_k b_k
  = \sum\limits_{k = 1}^{n - 1} (a_k - a_{k+1}) b_k + a_n B_n
  \label{th151:sum1}
  \end{gather}
\end{theorem}

\begin{proof}
  Очевидно, $b_1 = B_1, b_k = B_k - B_{k - 1}, k = 2, 3, \dots, n$ \\
  Поэтому
  $a_1 b_1 + a_2 b_2 + \dots + a_n b_n
  = a_1 B_1 + a_2 (B_2 - B_1) + a_3 (B_3 - B_2) + \dots a_n (B_n - B_{n - 1})
  = (a_1 - a_2) B_1 + (a_2 - a_3) B_2 + \dots + (a_{n-1} - a_n) B_{n-1} + a_n B_n$ \\
  Называется преобразованием Абеля $\sum\limits_{k = 1}^{n} a_k b_k$.
\end{proof}

\begin{consequence}[лемма Абеля]
  \label{th151:cons}
  Если $a_1 \leq a_2 \leq \dots \leq a_n$
   или $a_1 \geq a_2 \geq \dots \geq a_n$ \\
  $a_k \in \mathbb{R}, \ \forall k = 1, 2, \dots, n$,
  $|b_1 + \dots + b_k| \leq B, \ (b_k \in \mathbb{C})$, то
  $$\left|\sum\limits_{k = 1}^{n} a_k b_k \right| \leq B(|a_1| + 2|a_n|)$$
\end{consequence}
\begin{proof}
  $\left|\sum\limits_{k = 1}^{n} a_k b_k \right | \leq
  \sum\limits_{k = 1}^{n - 1} |a_k - a_{k + 1}| |B_k| + |a_n B_n| \leq
  B\left(\sum\limits_{k = 1}^{n - 1} |a_k - a_{k + 1}| + |a_n|\right) = \\
= B\left(\left| \sum\limits_{k = 1}^{n - 1} (a_k - a_{k + 1}) \right | + |a_n| \right)
= B(|a_1 - a_n| + |a_n|) \leq B(|a_1| + 2|a_n|).$
\end{proof}

\section{Признаки Дирихле и Абеля}
\begin{theorem}[признак Дирихле]
  \label{th161}
  Пусть дан ряд
  \begin{gather}
    \sum\limits_{n = 1}^{\infty} a_n b_n \label{th161:series1}
  \end{gather}
  \begin{enumerate}
    \item $a_n \in \mathbb{R}^n, b_n \in \mathbb{C}, \ n = 1, 2, \dots$
    \item $\{a_n\}, \{a_n\} \downarrow 0 \ (\{a_n\} \uparrow 0)$
    \item $\{B_n\}$ --- последовательность частичных сумм ряда $\sum b_n$ ограничена
  \end{enumerate}
  Тогда ряд \eqref{th161:series1} сходится.
\end{theorem}

\begin{proof}
  $\exists B > 0, \ |B_n| \leq B \ \forall n \Rightarrow \forall m \geq n \geq 2:
  |b_n + \dots + b_m| = |B_m - B_{n - 1}| \leq 2B$ \\
  Возьмем $\varepsilon > 0$. По скольку $a_n \to 0$, то
  $\exists N = N(\varepsilon): \forall n > N(\varepsilon)$ имеем
  $|a_n| < \frac{\varepsilon}{6B}$.\\
  Поэтому, $\forall n > N(\varepsilon)$ и $\forall m \geq n$ получим:
  $$|a_n b_n + \dots + a_m b_m| \leq 2B(|a_n| + 2|a_m|) <
  2B\left(\frac{\varepsilon}{6B} + 2 \frac{\varepsilon}{6B} \right) = \varepsilon$$
  Ряд \eqref{th161:series1} удовлетворяет Критерию Коши сходимости рядов.
\end{proof}
\begin{remark}
  Признак Лейбница - это частный случай признака Дирихле. \\
\end{remark}

\begin{theorem}[признак Абеля]
  \label{th162}
  Если последовательность действительных чисел $a_n$ монотонна и ограничена,
  ряд $\sum\limits_{n = 1}^{\infty} b_n, b_n \in \mathbb{C}$ сходится, то ряд
  \eqref{th161:series1} также сходится.
\end{theorem}

\begin{proof}
  $a_n = a + \alpha_n, \{\alpha_n\}$ --- монотонно стремящаяся к нулю
  последовательность. \\
  Поэтому
  $$\sum a_n b_n = \sum (a + \alpha_n)b_n = a \sum b_n + \sum \alpha_n b_n, $$
  где $a \sum b_n$ сходится по условию, а $\sum \alpha_n b_n$ сходится по
  признаку Дирихле. \\
  $\{B_n\}$ --- последовательность частичных сумм $\sum b_n$ ограничена,
  $\{\alpha_n\}$ --- монотонно стремящаяся к нулю последовательность.
\end{proof}

\section{Безусловно и условно сходящиеся ряды}
\begin{definition}
  Пусть $\{k_n\}, n = 1, 2, \dots$ --- последовательность, в которой каждое
  натуральное число встречается только один раз. $\{k_n\}$ --- однозначное
  отображение $a_n^* = a_{k_n}, \ (n~=~1, 2, \dots)$. \\
  Будем говорить, что ряд $\sum a_n^*$ является перестановкой ряда $\sum a_n$.
\end{definition}
\begin{definition}
  Говорят, что $\sum a_n$ сходится безусловно, если каждая перестановка
  сходится.
\end{definition}
\begin{theorem}
  \label{th171}
  Ряд $\sum a_n, (a_n \in \mathbb{C})$ сходится безусловно тогда и только тогда,
  когда он сходится абсолютно.
\end{theorem}

\begin{proof}
  Достаточность. \\
  Если ряд $\sum a_n$ сходится абсолютно, то все его перестановки сходятся к
  одному и тому же числу --- сумме исходного ряда. \\
  Пусть $\sum a_n^*$ --- перестановка ряда $\sum a_n$. $S_n^*$ --- ее частичная
  сумма. \\
  По Коши: $\forall \varepsilon > 0 \exists N : m \geq n > N$
  \begin{gather}
    |a_n| + \dots + |a_m| < \varepsilon \label{th171:uneq1}
  \end{gather}
  Выберем $p$ так, чтобы все натуральные числа $1, 2, \dots, N$ содержались в
  множестве $k_1, k_2, \dots, k_p$ (смотри определение), тогда при $n > p$
  $a_1, \dots, a_N$ в разности $S_n - S_n^*$ уничтожаются, так что
  $|S_n - S_n^*| < \varepsilon$ в силу \eqref{th171:uneq1}. \\
  Значит $\{S_n^*\}$ сходится к тому же пределу, что и $\{S_n\}$.
\end{proof}

\begin{definition}
  Сходящийся, но не абсолютно сходящийся ряд называется условно сходящимся.
\end{definition}
Из теоремы \eqref{th171} (из необходимости условия) $\Rightarrow$ Теорема \eqref{th172}
\begin{theorem}
  \label{th172}
  Условно сходящийся ряд не может сходится безусловно, то есть у него всегда
  существует расходящаяся перестановка.
\end{theorem}
\begin{proof}
  Без Доказательства.
\end{proof}

\begin{theorem}[Римана]
  Если ряд с действительными членами условно сходится, то каким бы не было
  действительное число $S$, существует перестановка ряда такая, что ее сумма
  равна $S$
\end{theorem}

\begin{proof}
  Без Доказательства.
\end{proof}

  \chapter{Функциональные последовательности и ряды}
  \section{Поточечная сходимость}
Пусть на некотором множестве $\mathbb{E}$ задана последовательность
комплексно значимых функций $f_n, n = 1, 2, \dots \ , \ (f_n \in \mathbb{C})$.
Элементы $x \in \mathbb{E}$ будем называть точками.
\begin{definition}
  $\{f_n\}$ называется ограниченной на $\mathbb{E}$, если
  $\exists M > 0: \forall n \in \mathbb{N}, \forall x \in \mathbb{E}$ выполняется
  $$|f_n(x)| \leq M$$
\end{definition}

\begin{definition}
  $\{f_n\}$ называется сходящейся поточечно на множестве $\mathbb{E}$, если
  при любом фиксированном $x \in \mathbb{E}$, числовая последовательность
  $\{f_n(x)\}$ сходится. Если последовательность сходится на $\mathbb{E}$,
  то $f(x) := \lim\limits_{n \to \infty} f_n(x), x \in \mathbb{E}$ называется
  пределом последовательности. Пусть
  $\{U_n(x)\}_{n = 1}^\infty, x \in \mathbb{E}, \ (U_n \in \mathbb{C})$
  --- последовательность числовых функций.
\end{definition}

\begin{definition}
  Множество числовых рядов
  \begin{gather}
    \sum\limits_{n = 1}^{\infty} U_n(x) \label{def213:series1}
  \end{gather}
  в каждой из которых точка $x$ фиксированная называется рядом на множестве
  $\mathbb{E}$, а функция $U_n(x)$ --- его член. \\
  $S_n(x) = \sum\limits_{k = 1}^{n} U_k(x), x \in \mathbb{E}$ называется
  $n$-ой частичной суммой ряда \eqref{def213:series1}. \\
  $\sum\limits_{k = n + 1}^{\infty} U_k(x)$ - его $n$-ым остатком.
\end{definition}

\begin{definition}
  Ряд \eqref{def213:series1} называется сходящимся поточечно на множестве
  $\mathbb{E}$, если последовательность $\{S_n(x)\}$ сходится поточечно на
  $\mathbb{E}$. При этом $\lim\limits_{n \to \infty} S_n(x) = S(x), x \in \mathbb{E}$
  называется суммой ряда \eqref{def213:series1}.
  $$S(x) = \sum\limits_{n = 1}^{\infty} U_n(x).$$
\end{definition}

\begin{definition}
  Если ряд \eqref{def213:series1} при любом $x \in \mathbb{E}$ сходится абсолютно,
  то он называется абсолютно сходящимся на множестве $\mathbb{E}$.
\end{definition}

\begin{remark}
  Беззаботная перестановка членов ряда может привести к ошибке.
\end{remark}

\section{Равномерная сходимость}
\begin{definition}
  \label{def221}
  Говорят, что функциональная последовательность $\{f_n\}_{n=1}^\infty$
  сходится равномерно на $\mathbb{E}$, если $\forall \varepsilon > 0 \
  \exists N \in \mathbb{N}: \forall n > N, \forall x \in \mathbb{E}$
  имеем
  $$|f_n(x) - f(x)| < \varepsilon$$
  Ясно, что каждая равномерно сходящаяся последовательность, сходится поточечно.
\end{definition}

\begin{comment}
  Обозначение равномерной сходимости:
  $f_n \stackrel{\mathrm{\mathbb{E}}}{\rightrightarrows} f$
\end{comment}

\begin{theorem}[Критерий Коши равномерной сходимости последовательностей]
  \label{th221}
  Для того, чтобы $\{f_n\}$ равномерно сходилась на
  $\mathbb{E} \Longleftrightarrow$
  $\forall \varepsilon > 0 \ \exists N: n, m > N, \forall x \in \mathbb{E}:$
  \begin{gather}
    |f_n(x) - f_m(x)| < \varepsilon \label{th221:uneq1}
  \end{gather}
\end{theorem}

\begin{proof}
  \hfill
  \begin{itemize}
    \item Необходимость: \\
      $f_n \stackrel{\mathrm{\mathbb{E}}}{\rightrightarrows} f$, тогда
      $\forall \varepsilon > 0, \ \exists N \in \mathbb{N}: \forall n > N,
      \forall x \in \mathbb{E} \ |f_n(x) - f(x)| < \frac{\varepsilon}{2}$. \\
      $|f_n(x) - f_m(x)| \leq |f_n(x) - f(x)| + |f(x) - f_m(x)| <
      \frac{\varepsilon}{2} + \frac{\varepsilon}{2} = \varepsilon$,
      $(\forall n, m > N, \forall x \in \mathbb{E}).$
    \item Достаточность: \\
      Пусть выполняется условие Коши, тогда $\{f_n(x)\}$, удовлетворяет
      критерию Коши сходимости числовых последовательностей и следовательно
      сходящегося числового предела, который обозначим $f(x)$. \\
      Тогда перейдя к пределу при $m \to \infty$ получим $\forall n > N,
      \forall x \in \mathbb{E}: |f_n(x) - f(x)| < \varepsilon$. \\
  \end{itemize}
\end{proof}
Иногда полезен критерий, следующий из определения \eqref{def221}

\begin{theorem}
  Пусть $\lim\limits_{n \to \infty} f_n(x) = f(x), \forall x \in \mathbb{E}$. \\
  Положим $r_n = \sup|f_n(x) - f(x)|, x \in \mathbb{E}$ --- равномерное уклонение. \\
  Тогда $f_n \stackrel{\mathrm{\mathbb{E}}}{\rightrightarrows} f
  \Longleftrightarrow r_n \to 0, \ n \to \infty$. (Переформулировка определения).
\end{theorem}

\begin{proof}
  Без доказательства. \\
\end{proof}

\begin{example}
  $f_n(x) = x^n, \mathbb{E} = [0, 1)$ \\
  $\lim\limits_{n \to \infty} f_n(x) = 0, \forall \in \mathbb{E},
  r_n = \sup\limits_{x \in [0, 1)} |x^n - 0| = 1 \not \to 0, n \to \infty.$ \\
  $\{x^n\}$ не является равномерно сходящейся на $\mathbb{E}$.
\end{example}

\begin{example}
  $f_n(x) = x^n - x^{n+1}, \mathbb{E} = [0, 1]$. \\
  $f_n(x) \to 0, \forall x \in \mathbb{E}, \  f_n'(x) = nx^{n-1} - (n + 1)x^n = 0$. \\
  $x_n = \frac{n}{n+1}, \ f_n(x_n) = x_n^n(1 - x_n) < \frac{1}{n + 1}$. \\
  $r_n < \frac{1}{n + 1}$. \\
\end{example}

\begin{definition}
  \label{def222}
  \begin{gather}
    \sum\limits_{n = 1}^{\infty} U_n(x), \ x \in \mathbb{E} \label{def222:series1}
  \end{gather}
  называется равномерно сходящейся, если на множестве $\mathbb{E}$ равномерно
  сходится последовательность частичных сумм. \\
\end{definition}

Пусть $S_k(x)$ --- частичные $k$-ые суммы ряда \eqref{def222:series1},
$$m \geq n: U_n(x) + \dots + U_m(x) = S_m(x) - S_n(x)$$
тогда из теоремы \eqref{th221} (критерий Коши равномерной сходимости последовательности)
$\Rightarrow$ Теорема \eqref{th223} (критерий Коши равномерной сходимости ряда).

\begin{theorem}[Критерий Коши равномерной сходимости ряда]
  \label{th223}
  Для того, чтобы ряд \eqref{def222:series1} равномерно сходился на множестве
  $\mathbb{E} \Longleftrightarrow \forall \varepsilon > 0 \ \exists N \in \mathbb{N},
  \forall n,m > N,
  \forall x \in \mathbb{E}: $
  \begin{gather}
    |U_n(x) + \dots + U_m(x)| < \varepsilon \label{th223:uneq1}
  \end{gather}
\end{theorem}

\begin{proof}
  Без доказательства.
\end{proof}

\begin{consequence}[Необходимый признак равномерной сходимости]
  У равномерно сходящегося ряда общий член равномерно стремится к нулю.
\end{consequence}

\begin{theorem}[Признак Вейерштрасса]
  Пусть $\{U_n\}$ --- последовательнсоть функций, определенных на $\mathbb{E}$
  и пусть $|U_n(x)| \leq a_n, \forall x \in \mathbb{E}, \forall n \in \mathbb{N}.$
  Тогда если $\sum a_n < \infty$ сходится, то следовательно $\sum U_n(x)$ сходится
  равномерно на $\mathbb{E}$.
\end{theorem}
\begin{proof}
  Если $\sum a_n$ сходится, то $\forall \varepsilon > 0 |
  \sum\limits_{k = n}^{m} U_k(x)| \leq
  \sum\limits_{k = n}^{m} a_k < \varepsilon$, при любом $x \in \mathbb{E}$,
  если только $m$ и $n$ достаточно велики, теорема \eqref{th111} (критерий Коши
  сходимости числового ряда). Равномерная сходимость нашего ряда вытекает из
  теоремы \eqref{th223}.
\end{proof}
\begin{remark}
  $\sum a_n$ называется мажорирующим рядом $\sum U_n(x)$.
\end{remark}

\begin{remark}
  ПРОВЕРИТЬ!!! \\
  Условие признака Вейерштрасса не являются необходимыми для равномерной
  сходимости ряда.
\end{remark}

\section{Признаки равномерной сходимости рядов Дирихле и Абеля}
\begin{theorem}
  Пусть дан ряд
  \begin{gather}
    \sum\limits_{n = 1}^{\infty} a_n(x) b_n(x), \ x \in \mathbb{E} \label{th231:series1}
  \end{gather}
  такой что:
  \begin{enumerate}
    \item $a_n(x) \in \mathbb{R}, \ b_n(x) \in \mathbb{C}, \ n = 1, 2, \dots$
    \item $a_n(x) \stackrel{\mathrm{\mathbb{E}}}{\rightrightarrows} 0$
      (Равномерная сходимость к нулю), $\{a_n(x)\}$ - монотонна.
    \item Последовательность частичных сумм $\{B_n(x)\} ряда \ \sum b_n(x)$ ограничена на множестве $\mathbb{E}$.
  \end{enumerate}
  Тогда ряд \eqref{th231:series1} равномерно сходится на множестве $\mathbb{E}$.
\end{theorem}

\begin{proof}
  В силу условия 3, $\exists B > 0: |B_n(x)| \leq B, \ \forall x \in \mathbb{E}, \
  \forall n \in \mathbb{N}$. \\
  $\forall x \in \mathbb{E}, m \geq n \geq 2: |b_n(x) + \dots + b_m(x)| =
  |B_m(x) - B_{n-1}(x)| \leq 2B$. \\
  $\forall \varepsilon > 0$ из условия 2 $\Rightarrow \exists N = N(\varepsilon):
  n > N(\varepsilon), \forall x \in \mathbb{E}$ выполняется неравенство:
  $$0 \leq |a_n(x)| < \frac{\varepsilon}{6B}.$$
  Примениев лемму Абеля \eqref{th151:cons}, получим:
  \begin{gather*}
    |a_n(x) b_n(x) + \dots + a_m(x) b_m(x)| \leq 2B(|a_n(x) + 2a_m(x)|) < \varepsilon \\ 
    \forall x \in \mathbb{E},
    m \geq n \geq N(\varepsilon)
  \end{gather*}
  В силу критерия Коши \eqref{th223}, ряд \eqref{th231:series1} сходится равномерно.
\end{proof}

\begin{theorem}[Признак Абеля]
  \label{th232}
  \begin{gather}
    \sum\limits_{n = 1}^{\infty} a_n(x) b_n(x) \label{th232:series1}
  \end{gather}
  \begin{enumerate}
    \item Если $a_n(x) \in \mathbb{R}, b_n(x) \in \mathbb{C}, n = 1, 2, \dots,
      x \in \mathbb{E}$.
    \item $\{a_n(x)\}$ ограничена на множестве $\mathbb{E}$ и монотонна
      $\forall x \in \mathbb{E}$.
    \item Ряд $\sum b_n(x)$ равномерно сходится на $\mathbb{E}$.
  \end{enumerate}
  Тогда ряд \eqref{th232:series1} равномерно сходится.
\end{theorem}

\begin{proof}
  Доказательство легко провести так, как была доказана теорема \eqref{th161}.
\end{proof}

\begin{example}
  $\sum\limits_{n = 1}^{\infty} \frac{c_n}{n^x}$ ---
  ряд Дирихле. \\
  Если этот ряд сходится в точке $x_0$, то он сходится равномерно
  $\forall x \in \mathbb{E}, \ \mathbb{E} = [x_0, +\infty)$. \\
  Можно воспользоваться Признаком Абеля:
  $$a_n(x) = \frac{1}{n^{x-x_0}}, \ b_n = \frac{c_n}{n^{x_0}}$$
\end{example}

\begin{exercise}
  Рассмотреть и доказать абсолютную сходимость при $x > x_0 + 1$
\end{exercise}

\section{Равномерная сходимость и непрерывность}

\begin{theorem}
  Пусть $f_n \stackrel{\mathrm{\mathbb{E}}}{\rightrightarrows} f, \ x_0$ ---
  предельная точка множества $\mathbb{E}$ и пусть
  $\lim\limits_{x \to x_0} f_n(x) = A_n, \ (n = 1, 2, \dots).$
  Тогда $\{A_n\}$ сходится и
  \begin{gather}
    \lim\limits_{x \to x_0} f(x) = \lim\limits_{n \to \infty} A_n \label{th241:lim1}
  \end{gather}
  Иными словами, 2 предельных перехода в данном случае коммутируют. \\
  $$\lim\limits_{x \to x_0} \lim\limits_{n \to \infty} f_n(x) =
  \lim\limits_{n \to \infty} \lim\limits_{x \to x_0} f_n(x)$$
\end{theorem}

\begin{proof}
  Пусть $\varepsilon > 0.$ В силу равномерной сходимости последовательности
  $\{f_n\} \ \exists N : n > N, m > N, x \in \mathbb{E},$
  \begin{gather}
    \left|f_n(x) - f_m(x)\right| < \varepsilon \label{th241:uneq1}
  \end{gather}
  Переходя в неравенстве \eqref{th241:uneq1} к приделу при $x \to x_0$ получим
  \begin{gather}
    \left|A_n - A_m \right| < \varepsilon, \ (n, m > N) \label{th241:uneq2}
  \end{gather}
  Поэтому $\{A_n\}$ --- последовательность для которой выполняется признак
  Коши сходимости последовательности $\Rightarrow$ она сходится. \\
  Обазначим ее предел $A$ \\
  \begin{gather}
    \left|f(x) - A\right| \leq \left|f(x) - f_n(x)\right| +
    \left|f_n(x) - A_n\right| + \left|A_n - A\right| \label{th241:uneq3}
  \end{gather}
  Выберем $n:$
  \begin{gather}
    |f(x) - f_n(x)| < \frac{\varepsilon}{3}, \ \forall x \in \mathbb{E}
    \label{th241:uneq4}
  \end{gather}
  Это возможно в силу равномерной сходимости.
  \begin{gather}
    |A_n - A| < \frac{\varepsilon}{3} \label{th241:uneq5}
  \end{gather}
  Затем, для этого $n$ подберем такую окрестность
  $U(x_0): x \in U(x_0), x \not = x_0,$ следовательно:
  \begin{gather}
    |f_n(x) - A_n| < \frac{\varepsilon}{3} \label{th241:uneq6}
  \end{gather}
  Из неравенств \eqref{th241:uneq3} --- \eqref{th241:uneq6} получим
  $$|f(x) - A| < \varepsilon, \ \forall x \in U(x_0), \ x \not = x_0$$
  Это равносильно равенству \eqref{th241:lim1}
\end{proof}

\begin{theorem}
  \label{th242}
  Последовательность функций, непрерывных в точке $x \in \mathbb{E}
  f_n \stackrel{\mathrm{\mathbb{E}}}{\rightrightarrows} f,$ то функция
  $f$ непрерывна в точке $x_0$.
\end{theorem}

\begin{proof}
  Без доказательства.
\end{proof}

\begin{remark}
  Обратное не верно, то есть последовательность непрерывных функций может
  неравномерно сходиться.
\end{remark}

Из теоремы \eqref{th242} и определения \eqref{def222} $\Rightarrow$ теорема \eqref{th243}

\begin{theorem}
  \label{th243}
  Если функции $U_n(x), \ (n = 1, 2, \dots), \ x \in \mathbb{E}$ непрерывны
  в точке $x_0 \in \mathbb{E}$ и ряд $\sum\limits_{n = 1}^{\infty} U_n(x)$
  равномерно сходится на $\mathbb{E}$, то его сумма $f(x)$ также непрерывна
  в точкке $x_0$.
\end{theorem}

\section{Равномерная сходимость и интегрирование}
\begin{theorem}
  \label{th251}
  Пусть $f_n$ --- последовательность действительных, значимых, интегрируемых на
  отрезке $[a, b]$ функций. Тогда функция $f$ также интегрируема на $[a,b]$ и
  \begin{gather}
    \int\limits_a^b f(x) dx = \lim\limits_{n \to \infty} \int\limits_a^b f_n(x) dx
    \label{th251:int1}
  \end{gather}
  Существование предела заранее не предполагается.
\end{theorem}

\begin{proof}
  $\forall \varepsilon > 0, \ \exists n:$
  \begin{gather}
    |f_n(x) - f(x)| < \varepsilon, \ x \in [a,b] \label{th251:uneq1}
  \end{gather}
  Зафиксируем $n$ и выберем разбиение
  $[a,b], \ \triangle_1, \dots, \triangle_S$ так, чтобы выполнялось неравенство
  \begin{gather}
    \sum_i \omega(f_n, \triangle_i) |\triangle_i| < \varepsilon \label{th251:sum1}
  \end{gather}
  \begin{comment}
    $\omega(f, E) = \sup - \inf$ --- колебание функции.
  \end{comment}
  Функции $f_n$ интегрируемы на $[a,b]$. По скольку
  $\omega(f, \triangle_i) \leq \omega(f_n, \triangle_i) + 2 \varepsilon,  \ (i = 1, \dots, S)$
  (смотри \eqref{th251:uneq1}).
  $$\sum_i \omega(f, \triangle_i)|\triangle_i| \leq \varepsilon + 2\varepsilon (b - a)$$
  Отсюда следует, что $f \in \mathbb{R} [a,b]$. Для доказательства \eqref{th251:int1}
  выберем $n > N: $
  \begin{gather*}
    |f_n(x) - f(x)| < \varepsilon, \ (a \leq x \leq b), \ n > N \\
    \left|\int\limits_a^b f(x) dx - \int\limits_a^b f_n(x) dx\right| \leq
    \int\limits_a^b |f(x) - f_n(x)| dx < \varepsilon(b - a)
  \end{gather*}
  Отсюда вытекает \eqref{th251:int1}. \\
\end{proof}

\begin{theorem}
  $U_n \in R[a,b]$ (Интегрируема). Если
  \begin{gather}
    f(x) = \sum\limits_{n = 1}^{\infty} U_n(x), \ (a \leq x \leq b)
    \label{th252:series1}
  \end{gather}
  При чем ряд \eqref{th252:series1} сходится на $[a,b]$, тогда
  $$\int\limits_a^b f(x) dx = \sum\limits_{n = 1}^{\infty} \int\limits_a^b
  f(x) dx$$
  Иными словами ряд \eqref{th252:series1} можно интегрировать частями.
\end{theorem}

\begin{proof}
  Без доказательства.
\end{proof}

\begin{remark}
  При нарушении равномерности ряд, состоящий из интегрируемых функций может
  иметь интегрируемую сумму.
\end{remark}

\section{Равномерная сходимость и дифференцирование}
$f_n(x) = \frac{\sin nx}{\sqrt{n}}, x \in \mathbb{R}$ показывает, что из
равномерной сходимости последовательности функций не следует даже поточечная
сходимость последовательностей функций производных. \\
То есть нужны более сильные предположения, чтобы заключать, что $f_n' \to f_n$,
при $f_n \to f$.

\begin{theorem}
  \label{th261}
  Пусть $f_n(x) \to f(x), \ x \in [a,b], \ n \to \infty, \ f_n \in C[a,b], \
  (n = 1, 2, \dots)$. \\
  Если $\{f_n'(x)\}$ сходится равномерно на $[a,b]$, то $f_n(x)$ дифференцируема
  и $$f'(x) = \lim\limits_{n \to \infty} f_n'(x)$$
\end{theorem}

\begin{proof}
  Обозначим через $f^*$ предел последовательности $f_n'$. Ввиду теоремы \eqref{th242}
  $f^*$ непрерывна на $[a,b]$. \\
  Применим теорему \eqref{th251} к последовательости $\{f_n\}$ на промежутке
  $[a, x],$ где $x \in [a,b]$
  $$\int\limits_a^x f^*(t) dt = \lim \int\limits_a^x f'(t) dt =
  \lim\limits_{n \to \infty} (f_n(x) - f_n(a)) = f(x) - f(a)$$
  Так как интеграл слева ввиду непрерывности функции $f^*$ имеет производную
  равную $f',$ то ту же производную имеет и $f(x)$.
  $$f'(x) = f^*(x) = \lim\limits_{n \to \infty} f'(x), x \in [a,b]$$
\end{proof}

Перефразируем теорему \eqref{th261} с точки зрения рядов: \\
Пусть сходящийся ряд $\sum\limits_{n = 1}^{\infty} U_n(x) =: f(x), x \in [a,b]$
и пусть $U_n(x) \in C^1[a,b], \ (n = 1, 2, \dots)$.\\
Если ряд $\sum\limits_{n = 1}^{\infty} U_n'(x)$ сходится равномерно на $[a,b]$,
то сумма $f(x)$ дифференцируема, и \\
$f'(x) = \sum\limits_{n = 1}^{\infty} U_n'(x), x \in [a,b]$.



\end{document}

